\defmodule{BASISSI}

Same as {\tt BASISLR}, except that the values are
represented using the {\tt SuperInteger} type.

\bigskip\hrule
\code

DEFINITION MODULE BASISSI;

FROM SUPINT   IMPORT SuperInteger;
FROM CONFIG   IMPORT BDim;

CONST
   BSizeSI = 2;                        (* Used in BCreateScal.            *)
TYPE
   BScal   = SuperInteger;
   BVect   = ARRAY BDim OF BScal;
   BMatr   = ARRAY BDim OF BVect;

VAR
   BZero, BOne, BMinusOne : BScal;     (* 0, 1, and -1                    *)

PROCEDURE BCreateScal (VAR S : BScal);
\endcode
  \tab  Creates the {\tt SuperInteger S}.
   The initial memory space (in words)
   allocated for the digits of {\tt S} is 2 to the
   power {\tt BSizeSI} (currently 4 words).
   (See {\tt Create} in {\tt SUPINT}).\endtab
\code

PROCEDURE BDeleteScal (VAR S : BScal);
\endcode
  \tab  Deallocates the memory allocated for the {\tt SuperInteger S},
  at a previous call to {\tt BCreateScal}. \endtab
\code

PROCEDURE BCreateVect (VAR V : BVect; d1, d2 : BDim);
\endcode
  \tab  Creates the {\tt SuperInteger} entries for indices 
  {\tt d1}  to  {\tt d2} inclusive
  (i.e., extends {\tt V[0..d1-1]} to {\tt V[0..d2]}).
  If  {\tt d2} is smaller than  {\tt d1}  then does nothing. \endtab
\code

PROCEDURE BDeleteVect (VAR V : BVect; d1, d2 : BDim);
\endcode
  \tab  Deallocates the memory allocated to the {\tt SuperInteger}
  elements of {\tt V[d1..d2]}.
  \endtab
\code

PROCEDURE BCreateMatr (VAR M : BMatr; d1, d2 : BDim);
\endcode
  \tab  Creates the {\tt SuperInteger} entries for indices 
   {\tt d1}  to  {\tt d2} inclusive; i.e., extends the square matrix 
   {\tt M[0..d1-1, 0..d1-1]} to {\tt M[0..d2, 0..d2]}.
   If {\tt d1} $> 0$, creates the missing entries in the {\tt d1} 
   vectors already there, in addition to creating the {\tt d2 - d1 + 1}
   missing vectors.
\endtab
\code

PROCEDURE BDeleteMatr (VAR M : BMatr; d1, d2 : BDim);
\endcode
  \tab  Does exactly the opposite of {\tt BCreateMatr (M, d1, d2)}.
  \endtab
\code

PROCEDURE BReadScal (VAR S : BScal);

PROCEDURE BWriteScal (S : BScal;   nb : CARDINAL);
\endcode
  \tab  Writes the contents of {\tt S} to the current output.
        Uses {\tt nb} digits per line.  \endtab
\code

PROCEDURE BWriteVect (VAR V : BVect;  n : BDim;  nb : CARDINAL);
\endcode
  \tab  Writes contents of {\tt V} to current output.
        Uses {\tt nb} digits per line for each number.  \endtab
\code

PROCEDURE BWriteMatr (VAR M : BMatr;  n : BDim;  nb : CARDINAL);
\endcode
  \tab  Writes contents of {\tt M} to current output.
        Uses nb digits per line for each number.  \endtab
\code

PROCEDURE BProdScalVect (s : BScal;  VAR V : BVect;  d : BDim);
PROCEDURE BProdScal (VAR U, V : BVect;  n : BDim;  VAR S : BScal);
PROCEDURE BProdScalLR (VAR U, V : BVect;  n : BDim) : LONGREAL;
PROCEDURE BProdScalSI (VAR U, V : BVect; n : BDim; Res : SuperInteger);
\endcode
\code
PROCEDURE BL1NormLR (VAR V : BVect;  n : BDim) : LONGREAL;
PROCEDURE BL1NormSI (VAR V : BVect;  n : BDim;  S1 : SuperInteger);
PROCEDURE BL2NormLR (VAR V : BVect;  n : BDim) : LONGREAL;
PROCEDURE BSupNormLR (VAR V : BVect;  n : BDim) : LONGREAL;
PROCEDURE BSupNormSI (VAR V : BVect;  n : BDim;  S1 : SuperInteger);
\endcode
\code
PROCEDURE BModifVectLR (VAR U, V : BVect;  c : LONGREAL;  n : BDim);
PROCEDURE BModifVectLI (VAR U, V : BVect;  c : LONGINT;  n : BDim);
PROCEDURE BModifVect (VAR U, V : BVect;  c : BScal;  n : BDim);
\endcode
\code
PROCEDURE BChangeSign (VAR S : BScal);
PROCEDURE BChangeSignVect (VAR V : BVect;  n : BDim);
PROCEDURE BAffect (S : BScal;  VAR T : BScal);
PROCEDURE BExchange (VAR S, T : BScal);
\endcode
\code
PROCEDURE BEqual (S, T : BScal): BOOLEAN;
PROCEDURE BSmaller (S, T : BScal): BOOLEAN;
PROCEDURE BGreater (S, T : BScal): BOOLEAN;
\endcode
\code
PROCEDURE BAbsSizeOrd (S, T : BScal): INTEGER;
PROCEDURE BAbs (S : BScal; VAR T : BScal);
PROCEDURE BSetNeg (VAR S : BScal);
PROCEDURE BInc (VAR S : BScal; i : LONGINT);
PROCEDURE BAdd (S, T : BScal; VAR Res : BScal);
PROCEDURE BMultiply (S, T : BScal; VAR Res : BScal);
PROCEDURE BDivideRound (S, T : BScal; VAR Res : BScal);
PROCEDURE BCopyVect (VAR U, V : BVect;  n : BDim);
PROCEDURE BCopyMatr (VAR M1, M2 : BMatr;  n : BDim);
PROCEDURE BConvLR (S : BScal): LONGREAL;
PROCEDURE BConvSI (R : BScal; S : SuperInteger);
PROCEDURE LIToBScal (i : LONGINT; VAR S : BScal);
PROCEDURE BAffectIdentMatr (VAR M : BMatr;  n : BDim);
PROCEDURE BTrMatr (M : BMatr; VAR Mtr : BMatr; d : BDim);
PROCEDURE BCanonicalVect (VAR V : BVect; i : BDim);
PROCEDURE BRandomScal (VAR s : BScal; m : BScal);

PROCEDURE BEuclide (A, B : BScal;  VAR C, D, E, F, G : BScal);
\endcode
  \tab For given $A$ and $B$, returns $C$, $D$, $E$, $F$, $G$ such that: 
   \begin {eqnarray*}
     C  A + D  B  &=& G = \mbox{gcd} (A,B) \\
     E  A + F  B  &=& 0.
   \end {eqnarray*}
where {gcd}  is the greatest common denominator.
  \endtab
\code

END BASISSI.
\endcode
